%% ============================================================
%%  §4  Constraint-Aware Generation
%%  Status: formal model locked; Coq implementation in verification/
%% ============================================================

\section{Constraint-Aware Generation}
\label{sec:generation}

This section specifies a \emph{generator} for an \spg{}~\(G\) over a constraint
domain~\(D\): a procedure that enumerates inputs \(s \in \Sigma^\ast\) accepted
by the parser of \S\ref{sec:engine-parsing}, paired with the evidence tree the
parser would derive on \(s\). The generator is the formal reference against
which the prefix parser is differentially tested. It is implemented and
mechanized in the Coq development under \verb|verification/generation.v|.

\subsection{Design intent}
\label{sec:gen-intent}

A naive enumerator would unroll syntactic derivations to a bounded depth and
post-hoc filter by the constraint evaluator of
Definition~\ref{def:constraint-domain}. This wastes work: syntactic
derivations that are semantically dead (\(\mathsf{Lost}\)) at the root are
recognizable from local rule evaluations. We therefore lift the constraint
domain into the generator itself: each production is presented to the
generator as a \emph{constrained production}, whose child positions carry
domain-supplied restrictions referencing already-realized siblings.
Generation descends through these constrained productions and prunes branches
whose verdict is \(\mathsf{Lost}\) before recursing further.

We deliberately keep the generator \emph{generic} over the domain: a new
domain plugs in by supplying a small interface (\S\ref{sec:gen-domain}) and
discharging one local proof obligation (\S\ref{sec:gen-obligations}). The
framework proofs (\S\ref{sec:gen-correctness}) are independent of the
particular domain.

\subsection{Generator domain interface}
\label{sec:gen-domain}

A \emph{generator domain} extends a constraint domain
\(D = (\mathsf{Rules}, \mathsf{Closed}, \mathsf{Ctx}, \mathit{eval}, \oplus)\)
with a symbolic restriction language and three operations.

\begin{definition}[Generator domain]
\label{def:gen-domain}
A \emph{generator domain} over a constraint domain \(D\) is a tuple
\[
  \mathcal{G}_D \;=\; (\mathsf{Restr},\; \mathsf{project},\; \mathsf{step},\; \mathsf{check}_\boxtimes),
\]
where:
\begin{itemize}[nosep]
  \item \(\mathsf{Restr}\) is a decidable language of \emph{restrictions} on
    child evidence summaries. A restriction may reference sibling
    positions by index.
  \item \(\mathsf{project} \;:\;
    \mathsf{Rules} \times \mathsf{Ctx} \times N \times \mathbb{N} \;\to\;
    \mathsf{Restr}^\ast\)
    maps a rule, a context, a nonterminal \(n\), and an alternative index \(a\)
    to a vector of per-position restrictions, one per symbol in \(p_a^n\).
  \item \(\mathsf{step} \;:\;
    (T \cup N) \times \mathsf{Restr} \times \pi \times \mathbb{N} \;\to\;
    (\Sigma^\ast \times \mathsf{Evidence})^\ast\)
    enumerates the realized values at a single child position. The argument
    \(\pi\) is the partial vector of already-realized siblings (left of the
    dot). The natural-number argument \(k\) is the \emph{value budget}: a
    domain-specific bound on enumeration depth (regex length for terminal
    recognizers, type size for typing-domain evidence, and so on).
  \item \(\mathsf{check}_\boxtimes \;:\;
    \mathsf{Restr}^\ast \times (\Sigma^\ast \times \mathsf{Evidence})^\ast \;\to\;
    \mathbb{B}\)
    is an optional cross-position consistency check applied after a full
    production has been realized.
\end{itemize}
\end{definition}

The Coq formalization expresses \(\mathcal{G}_D\) as a module type
\verb|GenDomain| declaring \(D\) as a sub-module of type
\verb|ConstraintDomain|, and four parameters \verb|Restriction|,
\verb|project|, \verb|step|, \verb|check_tuple|.

\subsection{Constrained production}
\label{sec:gen-cstr-prod}

\begin{definition}[Constrained production]
\label{def:cstr-prod}
Let \(p_a^n = \alpha_1[b_1]\cdots\alpha_m[b_m] \in A(n)\). The
\emph{constrained production} of \(p_a^n\) under context \(\Gamma\) is the
pair
\[
  \widehat{p}^{\,n,\Gamma}_a
  \;=\;
  \bigl(p_a^n,\; (r_1,\ldots,r_m)\bigr),
  \qquad
  (r_1,\ldots,r_m) = \mathsf{project}(\mathcal{T}(n), \Gamma, n, a),
\]
where \(\mathcal{T}(n)\) is the semantic rule attached to \(n\)
(Definition~\ref{def:constraint-domain}). Nonterminals without a semantic
rule yield restrictions \(r_i = \mathsf{trivial}\) for every position.
\end{definition}

\paragraph{Functional restrictions.}
\label{par:gen-functional}
The framework correctness theorems below require the following property of
the domain. A restriction \(r_i\) is \emph{functional} when, given the
realized siblings \(\pi = ((v_1, e_1), \ldots, (v_{i-1}, e_{i-1}))\), the set
\(\mathsf{step}(\alpha_i, r_i, \pi, k)\) does not depend on the resolution
order chosen for positions \(j < i\). A generator domain is \emph{functional}
when all restrictions produced by \(\mathsf{project}\) are functional. Under
this assumption, canonical left-to-right resolution
(\S\ref{sec:gen-algorithm}) is sound and complete with respect to all
resolution orders.

\begin{remark}[Symmetric relations and the subtype caveat]
\label{rem:gen-symmetric}
Equality, membership in a finite context, and unification at a fixed type
variable are functional. Non-symmetric ordering predicates (such as
\(\tau \subseteq \tau'\)) are not in general functional: their valid-set at a
given position depends on which side of the order has already been fixed.
The current typing-domain instance does not use such predicates in any
shipped grammar; should they appear, the canonical-order proof would have to
be replaced by an order-confluence argument.
\end{remark}

\subsection{Evidence tree}
\label{sec:gen-evtree}

The generator returns a string paired with an evidence tree, the canonical
inhabitant the parser would build on that string.

\begin{definition}[Evidence tree]
\label{def:gen-evtree}
An \emph{evidence tree} \(\mathcal{E}\) is inductively
\[
  \mathcal{E} \;::=\;
  \mathsf{ETLeaf}(x,\, r)
  \;\mid\;
  \mathsf{ETNode}(n,\, a,\, e,\, [\mathcal{E}_1, \ldots, \mathcal{E}_m])
\]
where \(x \in \Sigma^\ast\) is a segment text and \(r\) is the residual
recognizer at the leaf (\S\ref{sec:gram-def}); and where \(n \in N\),
\(a \in \{1,\ldots,k_n\}\), \(e \in \mathsf{Evidence}\), and the children
\(\mathcal{E}_i\) match the symbols of \(p_a^n\) left-to-right.
\end{definition}

\begin{definition}[Yield and tree depth]
\label{def:gen-yield-depth}
The \emph{yield} \(\lfloor \mathcal{E} \rfloor \in \Sigma^\ast\) is the
concatenation of the texts of the terminal leaves in left-to-right order.
The \emph{tree depth} \(\delta(\mathcal{E})\) is defined by
\(\delta(\mathsf{ETLeaf}(\cdot,\cdot)) = 0\) and
\(\delta(\mathsf{ETNode}(\cdot,\cdot,\cdot,[\mathcal{E}_1,\ldots,\mathcal{E}_m])) =
1 + \max_i \delta(\mathcal{E}_i)\).
\end{definition}

\subsection{Generation algorithm}
\label{sec:gen-algorithm}

We define the generator \(\mathsf{gen}\) as a depth-and-budget-bounded
recursive function. We use a canonical left-to-right resolution order
(Remark~\ref{rem:gen-symmetric}).

\begin{definition}[Generator]
\label{def:gen}
Let \(G\) be an \spg{} over a generator domain \(\mathcal{G}_D\). For
\(n \in N\), context \(\Gamma \in \mathsf{Ctx}\), value budget
\(k \in \mathbb{N}\), and depth budget \(d \in \mathbb{N}\), the
\emph{generator}
\[
  \mathsf{gen}_G \;:\; N \times \mathsf{Ctx} \times \mathbb{N} \times \mathbb{N}
  \;\to\;
  \mathcal{P}_{\!\!\text{fin}}(\Sigma^\ast \times \mathsf{EvidenceTree})
\]
is defined by
\[
  \mathsf{gen}_G(n, \Gamma, k, 0) \;=\; \emptyset,
\]
\[
  \mathsf{gen}_G(n, \Gamma, k, d+1)
  \;=\;
  \bigcup_{a=1}^{k_n} \mathsf{resolve}(G, n, a, p_a^n, \mathsf{project}(\mathcal{T}(n), \Gamma, n, a), \Gamma, k, d),
\]
where \(\mathsf{resolve}\) walks the production positions left-to-right; at
position \(i\) with symbol \(\alpha_i\) and restriction \(r_i\) and realized
prefix \(\pi_{<i}\):
\begin{itemize}[nosep]
  \item if \(\alpha_i \in T\), iterate over
    \(\mathsf{step}(\alpha_i, r_i, \pi_{<i}, k)\) and record each candidate
    \((x, e)\) as the leaf \(\mathsf{ETLeaf}(x, \partial_x \alpha_i)\);
  \item if \(\alpha_i \in N\), let
    \(\Gamma_i = \Gamma \oplus \nabla_{i-1} \oplus \cdots \oplus \nabla_1\)
    be the context after composing the right-bound effects of the realized
    siblings (Definition~\ref{def:constraint-domain}), and recurse:
    \((x_i, \mathcal{E}_i) \in \mathsf{gen}_G(\alpha_i, \Gamma_i, k, d)\);
  \item after all \(m\) positions are realized, evaluate
    \(\mathit{eval\_rule}(\mathcal{T}(n), \Gamma, [e_1, \ldots, e_m]) = (v, e, \nabla)\);
    discard the branch if \(v = \mathsf{Lost}\), else if
    \(\mathsf{check}_\boxtimes((r_1,\ldots,r_m), ((x_1, e_1),\ldots,(x_m, e_m))) =
    \top\), emit
    \(\bigl(x_1 \cdots x_m,\; \mathsf{ETNode}(n, a, e, [\mathcal{E}_1,\ldots,\mathcal{E}_m])\bigr)\).
\end{itemize}
\end{definition}

\paragraph{Termination.}
\(\mathsf{gen}\) terminates by structural recursion on \(d\). Each
nonterminal child recursion strictly decreases \(d\); terminal-position
enumeration is bounded by \(k\) and the finiteness of regex residuals at
depth \(k\) (\S\ref{sec:gram-def}); within a production, \(\mathsf{resolve}\)
recurses on the symbol list, which is finite.

\paragraph{Pruning.}
The early-discard step at \(v = \mathsf{Lost}\) is the safe pruning of
Lemma~\ref{lem:safe-pruning}: a branch whose root evaluates to \(\mathsf{Lost}\)
has empty denotation, hence no extension can produce a semantically valid
string.

\subsection{Correctness theorems}
\label{sec:gen-correctness}

\begin{theorem}[Generator soundness]
\label{thm:gen-sound}
Let \(G\) be an \spg{} over a functional generator domain \(\mathcal{G}_D\).
For every \(n,\Gamma,k,d\) and every
\((s, \mathcal{E}) \in \mathsf{gen}_G(n, \Gamma, k, d)\):
\begin{enumerate}[nosep]
  \item \(\mathcal{E}\) is a derivation tree of \(G\) at root \(n\)
    (Definition~\ref{def:gen-evtree});
  \item \(\lfloor \mathcal{E} \rfloor = s\);
  \item \(\delta(\mathcal{E}) \le d\);
  \item the constraint-domain evaluation \(\mathit{eval}\) at the root of
    \(\mathcal{E}\) under \(\Gamma\) is not \(\mathsf{Lost}\).
\end{enumerate}
\end{theorem}

\begin{theorem}[Generator completeness]
\label{thm:gen-complete}
Let \(G\) be an \spg{} over a functional generator domain \(\mathcal{G}_D\).
For every \(n,\Gamma,d\) and every evidence tree \(\mathcal{E}\) with root
\(n\), \(\delta(\mathcal{E}) \le d\), root evaluation not \(\mathsf{Lost}\)
under \(\Gamma\), and every terminal leaf text of length at most \(k\) in
the residual measure of its terminal recognizer, we have
\[
  \bigl(\lfloor \mathcal{E} \rfloor,\, \mathcal{E}\bigr)
  \;\in\;
  \mathsf{gen}_G(n, \Gamma, k, d).
\]
\end{theorem}

The leaf-length hypothesis is the price of an explicit budget: without it
the generator would have to enumerate infinitely many regex witnesses at a
single position (\S\ref{sec:gen-domain}).

\subsection{Domain-level proof obligation}
\label{sec:gen-obligations}

The framework theorems \ref{thm:gen-sound}--\ref{thm:gen-complete} depend on
a single lemma each domain must discharge. Informally:
\(\mathsf{step}\) returns exactly the values consistent with its restriction
under the realized prefix.

\begin{definition}[\(\mathsf{step}\) soundness obligation]
\label{def:step-sound}
A generator domain \(\mathcal{G}_D\) is \emph{step-sound} when for every
symbol \(\alpha\), restriction \(r\), realized prefix \(\pi\), budget \(k\),
and \((v, e) \in \mathsf{step}(\alpha, r, \pi, k)\), the pair \((v, e)\)
satisfies the domain-specific semantic interpretation
\(\llbracket r \rrbracket_\pi\) of \(r\) under \(\pi\).
\end{definition}

This is stated as an opaque \verb|Axiom step_sound| in the framework module
type and discharged inside each domain instance with its concrete restriction
language. The corresponding completeness obligation requires that every
\((v, e) \in \llbracket r \rrbracket_\pi\) whose value-budget cost is at most
\(k\) is enumerated by \(\mathsf{step}(\alpha, r, \pi, k)\).

\subsection{Differential testing of the parser}
\label{sec:gen-diff-test}

The generator is the testing oracle for the implementation of
\S\ref{sec:engine-parsing}. The differential pipeline is:
\begin{enumerate}[nosep]
  \item Extract \(\mathsf{gen}_G\) from Coq to OCaml.
  \item Run on a sample of shipped grammars (STLC, IMP, FUN) at varying
    \((k, d)\), producing a finite set of pairs \((s_i, \mathcal{E}_i)\).
  \item For each \(s_i\), run the Rust prefix parser on every token-boundary
    prefix; check that the parser accepts every prefix and that the final
    evidence tree matches \(\mathcal{E}_i\) up to evidence equality.
  \item A planned Rust-side enumerator will additionally enumerate
    \((s_i', \mathcal{E}_i')\) pairs and check set-equality against the Coq
    output at fixed \((k, d)\).
\end{enumerate}
The Coq generator is the specification: any divergence between Coq output
and Rust behavior is treated as a Rust bug.

\subsection{Coq module layout}
\label{sec:gen-coq-layout}

The Coq development mirrors the formal structure of this section:
\begin{center}
\begin{tabular}{ll}
  \toprule
  Coq object & This document \\
  \midrule
  \verb|ConstraintDomain|  & Definition~\ref{def:constraint-domain} \\
  \verb|GenDomain|         & Definition~\ref{def:gen-domain} \\
  \verb|cstr_prod|         & Definition~\ref{def:cstr-prod} \\
  \verb|evidence_tree|     & Definition~\ref{def:gen-evtree} \\
  \verb|gen|               & Definition~\ref{def:gen} \\
  \verb|gen_sound|         & Theorem~\ref{thm:gen-sound} \\
  \verb|gen_complete|      & Theorem~\ref{thm:gen-complete} \\
  \verb|step_sound|        & Definition~\ref{def:step-sound} \\
  \bottomrule
\end{tabular}
\end{center}
